\section{Error handling in Python}

\subsection{Errors and Exceptions}
\textbf{Error handling} is one of the most important problems to solve when designing a program. What should you do when a piece of code fails? What does fail mean? Some examples:
\begin{itemize}
    \item Invalid input e.g. passing a path to a non existent file, or passing a string to a function for dividing numbers
    \item Valid output not found, e.g searching the position of the letter 'd' in the string 'elephant'
    \item Output cannot be find in a reasonable amount of time
    \item Runtime resource failures: network connection down, disk space ended\dots
\end{itemize}

Historically, one can follow one of two philosophies when handling errors:
\begin{itemize}
    \item Return some \textbf{error flag} (in different ways) to tell the user that something went wrong;
    \item Use \textit{exceptions}.
\end{itemize}

\subsection{Error flags and their limits}

A typical convention for programs is to return 0 from the main if the execution was successful and an error code (integer number) otherwise. However this isn't the best practice. Take a look a this example:

\input{pygments/error_flags}

Error codes have their use and thay are fine in some cases, but they suffer from a few issues:
\begin{itemize}
    \item Choosing them is often arbitrary and sometimes it is difficult to make a sensible choice. What if all the numbers can represent meaningful output of the function?
    \item They are cumbersome to use. Which error flag is used by a function? 0? -1? 99999999? You have to go through the documentation for each! If you have a deep hierarchy of functions you have to perform checks and pass the error up at every level!
    \item What if the caller of a function does not check the error flag? The bug can propagate \textit{silently} through its code!
\end{itemize}

\subsection{Exceptions}

With that in mind, we want something that
\begin{itemize}
    \item is clearly separated from the returned output
    \item cannot be silently ignored by the user
    \item is easy to report to upper level without lots of lines of code.
\end{itemize}
This is where \textit{exceptions} come into play:

\input{pygments/exceptions_vs_err_flags}

An \textbf{exception} is an object that can be \textbf{raised} (in other languages also \textit{thrown}) by a piece of code to signal that something went wrong. When an exception is raised, the normal flow of the code is interrupted. The program automatically propagates the exception back in the function hierarchy until it finds a place where the exception is \textbf{caught} and handled. If the exception is never caught, not even in the main, the program crashes \textbf{with a specific error message}. Catching the exception is done in Python with a \texttt{try - except} block, just like this:

\input{pygments/exceptions_brief}

There are two more optional statements in a \texttt{try-except} block:
\begin{itemize}
    \item \texttt{else}: contains a piece of code that is executed only if no exception is raised in the try block;
    \item \texttt{finally}: contains a piece of code that is executed no matter what.
\end{itemize}
\texttt{finally} is executed even if there is a return statement in the try block and can be used to release important resources (e.g. closing a file, or a connection).

\input{pygments/exceptions}

If that was all, exceptions would only be moderately useful. The real bargain is that \textit{you can send back information} together with the exception. In fact you \textit{are sending a full object}: the exception itself. Inside the exception you can report all kind of data useful to reconstruct the exact error, which can be used by the caller for debugging or to produce meaningful error messages. You can also select which exceptions you catch, leaving the others propagating up, as done in this example: 

\input{pygments/try_block}

Python provides a rich hierarchy of exception classes, which you can further customize (if you want) by deriving your own subclasses. Here's a schematic representation:
\begin{figure}[H]
\centering\includegraphics[width=\textwidth]{python_exceptions}
\end{figure}

However, even exceptions have some caveats. For example, catching \texttt{Exception}, will also catch \texttt{SyntaxError} and \texttt{NameError}. This means that the code will 'run' even if there is a typo in it! Bottom line: \textbf{you should never catch generically for \texttt{Exception}}, always be more specific.
Even worse, you should never catch for \texttt{BaseException} as that would even prevent the user from aborting the execution with a \texttt{KeyboardInterrupt} (e.g. Ctrl-C). 

In Python, exceptions are the default methods for handling failures. Many functions raise an exception when something goes wrong. The common approach is: do not check the input beforehand, use it and be ready to catch exceptions if any.

\input{pygments/dont_ask_permission}

\subsection{The traceback}
A \textbf{stack traceback}, or simply traceback, is a list of the function calls at a specific point in the code. An exception always stores the traceback at the failure point, since it contains crucial debug information.

In Python a traceback is read from the bottom up:
\begin{itemize}
    \item The last line of the traceback is the error message line, which contains the name of the exception that was raised;
    \item Further up are the various function calls from most recent to least recent.
\end{itemize}
When you catch an exception, the content of the traceback can be displayed using either the \texttt{logging} module or the \texttt{traceback} module. The following snippets of code will show how to work with these two modules:

\subsubsection{- Traceback report with \texttt{logging}}
\input{pygments/traceback_logging}

\subsubsection{- Traceback report with \texttt{traceback}}
\input{pygments/traceback_tb}

\subsection{- Exception logging with \texttt{traceback}}
\input{pygments/traceback_exception}

\subsection{Raising and customizing exceptions}

Up to now we have been dealing with exceptions generated by Python functions. What about raising exceptions ourselves?

\input{pygments/raising}

Besides the built-in exceptions provided by Python, you can add your own custom exceptions by inheriting from the \texttt{Exception} class. This serves two purposes:
\begin{itemize}
    \item Makes the exception handling code more specific, and hence more readable;
    \item Allows you to pass additional data with your exception - in the form of attributes of the class - which can be used for debugging or any other purpose.
\end{itemize}
Here's how you can create and use your own exception:
\input{pygments/custom_exceptions}
\hrulefill
\input{pygments/custom_exceptions_2}

One last question: when should you catch exceptions? Differently from error flags, which need to be checked as early as possible, you are not in a rush with exceptions. Remember: your goal is to provide the user a meaningful error message and useful debug information. You should catch an exception only when you have enough context to do that, which sometimes means waiting a few levels in the hierarchy! Here's a long example to let you understand this concept:

\begin{figure}[H]
\centering
\includegraphics[width=0.8\textwidth]{lab_file.png}
\end{figure}

\input{pygments/when_to_catch}

\hrulefill

\subsubsection{Catch too early}
\input{pygments/when_to_catch_1}

\hrulefill

\subsubsection{Catch when needed}
\input{pygments/when_to_catch_2}

\newpage 
\section{Generators}

\subsection{Iterators and iterables}
Let's analyse the properties of iterables a little deeper. Here's a little reminder of what we already know:
\begin{itemize}
    \item An \textbf{iterable} in Python is something that has a \texttt{\_\_iter\_\_} method, which returns an \textit{iterator}
    \item An \textbf{iterator} is an object that implements a \texttt{\_\_next\_\_} method which is used to retrieve elements one at a time
\end{itemize}
Here's more information about iterators:
\begin{itemize}
    \item When there are no more elements to return, the iterator signals that with a specific exception: \texttt{StopIteration()}
    \item An iterator also implements a \texttt{\_\_iter\_\_} method that returns itself. So an iterator is also technically an iterable\footnote{Only technically because an iterator has no data of its own, so you always need a "real" iterable to actually iterate}, but the opposite is not true.
\end{itemize}
Take a look at the following example where we loop over all the elements of a list (iterable) using \texttt{for} and \texttt{while}:

\input{pygments/show_iterator}

In the following examples we implemented two iterators of our own and tested their behaviour. Look at their structure and try and write one of your own:

\textbf{- A simple iterator}
\input{pygments/simple_iterator}
\hrulefill
\input{pygments/test_simple_iterator}

\bigskip

\textbf{- A crazy iterator}
\input{pygments/crazy_iterator}
\hrulefill
\input{pygments/test_crazy_iterator}

\subsection{Generators}

We have seen that iterators are useful to iterate over container. However that assumes a containers exists which means slots of memory are being used.
\textbf{Generators} allow you to loop over sequences of items even when they don't exist before - the items are just created \textbf{lazily} the moment they are required. For example you can write a generator to loop over the Fibonacci succession. You can't create the sequence earlier, since it is not finite!

Generators are created through either \textit{generator expressions} or \textit{generator functions}. In real life most of the time you will simply use pre-made functions that return a generator, like \texttt{range()} (in Python 3). Generators can be used to iterate in for loops, just like iterators. Take a look at these simple generator expressions:

\input{pygments/generators}

A \textbf{generator function} is a function that contains the keyword \texttt{yield} at least once in his body. When you call a generator function the code is not executed - instead a generator object is created and returned (even if you don't have a return statement). Each call to \texttt{next()} on the returned generator will make the function code run until it finds a \texttt{yield} statement. Then the execution is paused and the value of the expression on the right of \texttt{yield} is returned (yielded) to the caller. A further call of next will resume the execution from where it was suspended until the next \texttt{yield} and so on. Eventually, when the function body ends, \texttt{StopIteration} is raised. Usually generators functions contain a loop, but it's not mandatory!
Here's two examples of generator functions:

\textbf{- Simple generator function}
\input{pygments/generator_functions}

\textbf{- Infinite sequence generators}
\input{pygments/fibonacci}

Take a look at the list of functions in Sec. \ref{sec:special_iterable}: Many of those return a generator from an iterable, such as: \texttt{enumerate}, \texttt{map}, \texttt{filter}, \texttt{zip}, \texttt{reversed} and all the functions mentioned from the \texttt{itertools} library. Look at the documentation of each function if you want to see how to properly call it. In this next example we will showcase the capabilities of \texttt{itertools}:

\input{pygments/itertools_showcase}

\hrulefill
\subsubsection{Note: Lambda functions}

\textbf{Anonymous functions}, or \textbf{lambda functions}, are a construct typical of \textbf{functional programming} (see for example \url{https://en.wikipedia.org/wiki/Lambda_calculus} and \url{https://en.wikipedia.org/wiki/Functional_programming}).

In Python a lambda function is essentially a special syntax for creating a function on the fly, without giving it a name. They are limited to \textit{a single expression}, which is returned to the user. Many of the typical uses for lambdas are already covered in Python by generator expressions and comprehension, so this is more like a niche feature of the language. Here's an example:

\input{pygments/lambda}
\hrulefill

This next example simulates a programming session in which the user tries to create a more complex iterator, simplifying its structure step by step:

\textbf{- Recap example: file iterator}
\input{pygments/file_iterator}
\bigskip

\textbf{- File iterator redone}
\input{pygments/file_iterator_2}
\bigskip

\textbf{- File iterator, final version}
\input{pygments/file_iterator_3}